\chapter{State of the Art}



\section{Neuroevolution and Genetic Algorithms}

Genetic Algorithms (GAs) are search heuristics inspired by Darwinian evolution, utilizing mechanics such as population dynamics, fitness evaluation, crossover, and mutation. Neuroevolution (NE) is the specific application of GAs to Artificial Neural Networks (ANNs). It focuses on mimicking the evolution of a biological brain by optimizing either the synaptic weights or the overall network architecture to achieve a desired outcome.

Historically, a prominent approach in this field has been Topology and Weight Evolving Artificial Neural Networks (TWEANNs), which simultaneously evolve the structure and the connections of a network. However, traditional TWEANNs suffered from three critical limitations. First, they struggled heavily with the \textit{Competing Conventions} problem (also known as the permutation problem). This occurs when two parent networks compute the same logical solution but possess different internal topological orderings, crossing them over often produces dysfunctional offspring with scrambled or missing features. Second, early TWEANNs lacked a mechanism to preserve topological innovations. Without dividing the population into different species, newly mutated (and temporarily weaker) structures were quickly driven to extinction by older, highly optimized, but structurally simpler genomes. Lastly, to ensure genetic diversity, early models frequently initialized populations with complex, randomized topologies. This bloated the genomes with unjustified connections and unnecessarily expanded the search space.

The NeuroEvolution of Augmented Topologies (NEAT) algorithm \cite{stanley2002} introduces three core innovations to systematically resolve these issues: historical markings, speciation, and minimal initial structures. By assigning a global innovation number (an ID) to every new structural mutation, NEAT can perfectly align and cross over genomes of different sizes without requiring complex topological analysis, effectively eliminating the Competing Conventions problem. Furthermore, NEAT employs speciation (or niching) to group structurally similar networks together. This protects new, innovative topologies, giving them time to optimize their weights before they must compete with the entire population. Finally, NEAT addresses the initialization problem by starting the population with zero hidden nodes. Complexity is added incrementally only when it provides a distinct functional benefit, ensuring the final network remains as compact as possible and drastically reducing the dimensionality of the search space.

Through extensive ablation studies and benchmark testing on complex control tasks---such as the XOR problem and double pole balancing---NEAT has consistently demonstrated superior performance over traditional NE methods. Its adaptability has also been proven in dynamic video game environments, notably in \textit{Super Mario Bros.} \cite{sethbling2015mario}, which serves as the foundational inspiration for the baseline environment in this project. Despite its structural efficiency, however, standard NEAT possesses one major limitation: it relies on continuous mathematical activation functions (e.g., Sigmoid). These continuous floating-point operations are computationally expensive, making standard NEAT unsuitable for edge-computing scenarios where extreme energy efficiency is mandatory, such as robotics or aerospace applications. This computational bottleneck necessitates the transition toward Spiking Neural Networks.



\section{Spiking Neural Networks (SNN) and Green AI}

As previously mentioned, the world is currently experiencing the rapid rise and commercialization of AI and GPAI. However, there is ongoing dispute regarding the computational cost of running these machines and the vast volumes of data they require. Most of these modern architectures fall under the category of Red AI. According to Schwartz et al., ``Red AI refers to AI research that seeks to obtain state-of-the-art results in accuracy (or related measures) through the use of massive computational power---essentially `buying' stronger results'' \cite{schwartz2020green}. While Red AI focuses primarily on maximizing accuracy through sheer computational force, Green AI advocates for treating energy efficiency and computational cost as primary evaluation metrics alongside accuracy \cite{schwartz2020green}. Consequently, this project evaluates energy efficiency as a core metric, operating under the premise that it is just as valuable as model accuracy.

Furthermore, the primary reason these standard implementations are so computationally (and therefore energetically) expensive is the activation function. The activation function is a mathematical equation that determines whether a given node should remain inactive or fire, defining its output. The vast majority of Red AI alternatives utilize continuous activation functions, such as Sigmoid or ReLU. The fundamental issue is that these functions rely on continuous floating-point numbers, requiring the underlying hardware to perform complex matrix floating-point operations constantly.

As Schwartz et al. observe: ``Ironically, deep learning was inspired by the human brain, which is remarkably energy efficient'' \cite{schwartz2020green}. As mentioned previously, a biological brain does not physically operate like an open tap, rather, it processes information via discrete spikes. A biological neuron is divided into three main components: dendrites, the soma, and the axon. Roughly speaking, the dendrites can be considered the ``input system,'' the soma acts as the central processor or ``GPU,'' and the axon serves as the ``output system'' \cite{gerstner2014}. Electrical energy travels through these components, and a sufficient accumulation of voltage triggers an action potential. A group of these action potentials over a given time period constitutes a spike. These biological mechanics form the foundational principles of Spiking Neural Networks (SNNs) \cite{gerstner2014}. Utilizing spikes as an activation mechanism renders the architecture significantly more energy-efficient, as it relies on discrete, binary event-driven signals rather than continuous floating-point operations. Furthermore, SNNs are fundamentally inspired by electrical circuitry, making them highly optimized for neuromorphic hardware translation.

However, for an SNN to function computationally, it requires a physical mathematical abstraction. In this project, that abstraction is the Leaky Integrate-and-Fire (LIF) model \cite{gerstner2014}. As the name suggests, the LIF model governs neuronal dynamics through three distinct stages: integrate, leak, and fire. First, the neuron accumulates energy from incoming synaptic spikes, which raises its internal membrane potential (integration). Second, the membrane possesses a natural electrical resistance, if spikes do not arrive fast enough, the membrane progressively loses voltage, returning toward its resting state (leak). Finally, if the integration outpaces the leak and surpasses a strict voltage threshold, the neuron instantly emits a binary spike and resets its internal voltage to a baseline (fire). In this architecture, every continuous node generated by the NEAT algorithm is replaced by this discrete LIF physical model.

This section has outlined the significant computational advantages of SNNs and the LIF model, demonstrating why they serve as a powerful framework for achieving Green AI efficiency. While combining the NEAT algorithm with LIF neurons yields an exceptionally efficient Spiking-NEAT architecture, it introduces a critical challenge: this discrete, non-differentiable implementation cannot be trained using standard continuous mathematics, such as Backpropagation. A Spiking-NEAT architecture requires a biologically plausible learning rule to optimize its weights, which leads to the final core component of this project: biological Reinforcement Learning.



\section{Synaptic Plasticity and Biological Reinforcement Learning}

Most of the conventional AI architectures discussed previously rely on continuous activation functions and smooth, differentiable curves, therefore, the standard method for training these Artificial Neural Networks (ANNs) is Backpropagation. Because Backpropagation calculates gradients using loss and error functions, it strictly requires continuous values. In contrast, the Spiking-NEAT approach utilizes non-differentiable, discrete binary signals (spikes). Consequently, traditional gradient-descent methods cannot be applied. Furthermore, relying solely on NEAT's random weight mutations between evolutionary generations is computationally slow and inefficient. Therefore, this project proposes an alternative, biologically plausible solution to dynamically evolve the synaptic weights while the agent interacts with the environment.

In biological brains, learning is driven by synaptic plasticity, which refers to the ability of neural connections to strengthen or weaken over time based on their activity. This correlates with the Hebbian postulate: ``Neurons that fire together, wire together.'' This rule states that if neuron A repeatedly and persistently takes part in firing neuron B, the synaptic efficacy from A to B increases \cite{gerstner2014}. These principles are mathematically formalized in Spike-Timing-Dependent Plasticity (STDP). In STDP, the change in synaptic weight depends heavily on the precise timing of spikes. Generally, if a pre-synaptic neuron fires just before a post-synaptic neuron, the connection between them is strengthened, a process known as Long-Term Potentiation (LTP). Conversely, if the pre-synaptic neuron fires just after the post-synaptic neuron, the connection is weakened, resulting in Long-Term Depression (LTD).

However, standard STDP is strictly unsupervised, it evaluates only the relative timing of spikes, lacking any mechanism to determine whether the resulting action was beneficial or detrimental to the system. Consequently, if a negative action occurs repeatedly, STDP will continually strengthen the synaptic pathways responsible for that action simply because the neurons fire in sequence. Without a goal-oriented mechanism, standard STDP can reinforce maladaptive behaviors, ultimately hindering the agent's progression in dynamic environments.

To resolve this credit assignment problem, Reward-Modulated STDP (R-STDP) is introduced \cite{mozafari2018}. R-STDP integrates the principles of Reinforcement Learning by expanding standard STDP into a three factor learning rule. The three key factors are: the pre-synaptic spike, the post-synaptic spike, and a delayed neuromodulatory signal, which mathematically mimics the release of dopamine in biological brains. In this model, when neurons fire sequentially, STDP initially marks the connection with a temporary ``eligibility trace.'' If the agent subsequently performs a beneficial action and receives an environmental reward (a dopamine boost) shortly after, this trace is converted into a permanent weight increase (LTP). Conversely, if the agent receives a punishment or no reward, the connection is actively weakened (LTD) \cite{mozafari2018}. Consequently, R-STDP not only evolves the network weights dynamically but also ensures that learning is goal-directed and efficient.

In conclusion, this triad forms the complete theoretical foundation of the proposed model: NEAT acts as the architect to build the minimal network topology, LIF neurons serve as the highly energy-efficient processing engine using discrete spikes, and R-STDP functions as the biological reinforcement learning mechanism, adjusting synaptic weights dynamically based on behavioral success rather than isolated timing.



\section{Comparative Analysis and Gap Identification}

In summary, the theoretical frameworks discussed thus far establish several key premises: NEAT is a highly effective neuroevolutionary algorithm that dynamically evolves topologies but incurs high computational and energetic costs. Spiking Neural Networks, specifically utilizing the LIF model, offer a paradigm shift toward energy efficiency by replacing continuous floating-point operations with discrete binary spikes; however, these non-differentiable networks cannot be trained via Backpropagation. Finally, R-STDP provides a biologically plausible reinforcement learning mechanism to evolve synaptic weights dynamically, overcoming the limitations of unsupervised STDP by distinguishing between beneficial and detrimental actions.

The individual viability of these mechanisms has been validated in recent literature. Qiu et al. \cite{qiu2018evolving} demonstrated that a Spiking-NEAT architecture is not only possible but can outperform standard NEAT in specific control tasks. Similarly, Mozafari et al. \cite{mozafari2018} provided substantial evidence that R-STDP is an effective learning rule for SNNs. However, despite these isolated successes, a distinct research gap remains: these methodologies have rarely been unified and evaluated together within a highly complex, dynamic environment. Therefore, this project aims to prove that the integration of a Spiking-NEAT architecture, LIF neuron models, and R-STDP learning is a viable and highly efficient solution for complex control tasks.

To benchmark this integration, the agent is tasked with navigating \textit{Super Mario Bros.}, an environment directly inspired by the neuroevolution showcase MarI/O \cite{sethbling2015mario}. While utilizing a commercial video game may intuitively seem informal, such environments have become the gold standard for benchmarking modern reinforcement learning models. Unlike simple balancing tasks, \textit{Super Mario Bros.} is a highly dynamic, non-linear control environment characterized by a high-dimensional observation space. Because the agent must interpret the game state from raw visual pixel data, the computational burden of processing spatial relationships, momentum-based physics, and moving obstacles is exceptionally high. This complexity makes it an ideal testing ground to validate whether an energy-efficient Spiking Neural Network can match the performance of computationally expensive standard neural architectures in complex, real-time tasks. Furthermore, successfully solving this environment serves as a proof of concept for broader applications in robotics and aerospace, where extreme energy constraints and high-dimensional data processing are critical.

Ultimately, the objective of this thesis is to provide a thorough comparative analysis, demonstrating how a Spiking-NEAT architecture guided by R-STDP can achieve comparable or superior results to a standard NEAT baseline, while significantly reducing overall energy consumption and accelerating convergence.